\chapter{INTRODUÇÃO}
A Introdução apresenta uma visão geral do trabalho, permitindo ao leitor compreender o contexto, o problema investigado, a justificativa, os objetivos, a metodologia adotada e a organização da monografia. Embora seja o primeiro capítulo do trabalho, recomenda-se que sua revisão final seja feita após a conclusão do TCC.

Na parte inicial da Introdução (antes das seções específicas), recomenda-se elaborar parágrafos curtos e objetivos contemplando a seguinte sequência lógica:
\begin{uemaenum}
    \item \textbf{Contextualização e Tema:} Apresente o tema central da sua monografia, definindo os conceitos principais e situando o leitor sobre o cenário tecnológico ou acadêmico atual desse assunto;
    \item \textbf{Problematização:} Apresente a situação-problema encontrada (a lacuna ou dificuldade real observada) e culmine na sua pergunta norteadora (o problema de pesquisa);
    \item \textbf{Justificativa:} Explique a relevância do tema e a motivação do estudo. Por que esta pesquisa ou solução tecnológica é importante para a sociedade, para o mercado e para o meio acadêmico?
\end{uemaenum}

\section{Objetivos}

\subsection{Objetivo Geral}
Apresente aqui o objetivo principal do trabalho em um único parágrafo conciso, iniciado por um verbo no infinitivo (ex.: desenvolver, analisar, propor, implementar, avaliar), indicando claramente o propósito central da pesquisa ou do sistema proposto.

\subsection{Objetivos Específicos}
Os objetivos específicos detalham as etapas necessárias para alcançar o objetivo geral. Podem ser apresentados após uma breve frase introdutória, utilizando alíneas ou marcadores:
\begin{uemaball}
    \item levantar e analisar os requisitos funcionais e não funcionais do sistema;
    \item modelar a arquitetura e o banco de dados da aplicação;
    \item implementar o protótipo ou sistema utilizando as tecnologias definidas;
    \item avaliar e validar a solução desenvolvida junto aos usuários ou por meio de testes.
\end{uemaball}

\section{Metodologia}
Descreva de forma clara os procedimentos metodológicos adotados para a condução da pesquisa e o desenvolvimento do trabalho. Indique a natureza e a classificação da pesquisa (ex.: aplicada, exploratória, descritiva, bibliográfica, estudo de caso), bem como as etapas de desenvolvimento de software, levantamento de requisitos, ferramentas utilizadas e forma de validação ou coleta de dados.

\section{Estrutura do documento}
Este trabalho está estruturado da forma como segue. Além deste capítulo introdutório, no Capítulo 2 é apresentada a fundamentação teórica e os trabalhos correlatos. O Capítulo 3 descreve o desenvolvimento do trabalho, abrangendo a arquitetura, modelagem e implementação da solução. O Capítulo 4 apresenta os resultados obtidos e as discussões. Finalmente, no Capítulo 5 são apresentadas as conclusões e as propostas para trabalhos futuros.
